% !TEX root = ../main.tex
\chapter{Comparative Analysis of Quantum Technologies}
\label{ch:comparative_analysis}

{
\renewcommand{\abstractname}{Summary}
\renewenvironment{abstract}{%
  \par\noindent
  {\bfseries\large\abstractname}\par  % Titolo grassetto a sinistra
  \vspace{0.5em}
  \noindent
}{%
  \par\vspace{1em}
}
\begin{abstract}
The previous chapters have described the main pillars of this work—quantum sensing, quantum hardware, and quantum algorithms—each representing a distinct layer of the quantum technology stack. This wrap-up chapter provides a cross-domain comparative analysis integrating these perspectives to highlight synergies, scalability challenges, and performance trade-offs across physical implementations and computational paradigms. The goal is to establish a quantitative framework to evaluate the readiness and integration potential of different quantum technologies.
\end{abstract}
}
\section{Performance Metrics}
\subsection*{Hardware Metrics}
For quantum processors, key figures of merit include coherence times (\(T_1, T_2\)), gate and readout fidelities (\(F_g, F_r\)), operational temperature (\(T_{\mathrm{op}}\)), connectivity, and qubit frequency dispersion. Scalability depends on the ability to maintain fidelity under system growth and to mitigate cross-talk and control errors.

\subsection*{Algorithmic Metrics}
For quantum algorithms and hybrid machine learning architectures, performance metrics include circuit depth, number of qubits, training convergence, final accuracy, and robustness to noise. Hardware-aware benchmarks consider the depth-fidelity product and total execution time per iteration.

\subsection*{Sensing Metrics}
Quantum sensors are characterized by sensitivity (minimum detectable signal per unit bandwidth), dynamic range, and bandwidth. The quantum Fisher information (QFI) formalism enables technology-agnostic comparison between sensing modalities. A higher QFI implies improved precision for the same number of resources.



\section{Cross-Domain Quantitative Comparison}
Table~\ref{tab:hardware_compare} summarizes superconducting hardware properties, while Table~\ref{tab:algorithm_compare} presents computational metrics for algorithmic families. Lastly, table~\ref{tab:sensing_compare} compares representative quantum sensing platforms.

\begin{table}[h]
\centering
\caption{Superconducting qubit performance metrics~\cite{kim2021, place2021, kjaergaard2020, krantz2019}.}
\resizebox{\textwidth}{!}{
\begin{tabular}{lcccc}
\toprule
Material & $T_1$ [$\mu$s] & Gate Fidelity & $T_{\mathrm{op}}$ [K] & Typical Frequency [GHz] \\
\midrule
Al/AlO$_x$/Al (planar) & 100--200 & 0.999 & 0.02 & 5--7 \\
Ta-based transmon & 200--400 & 0.9995 & 0.02 & 4--6 \\
NbN/AlN/NbN (this work) & 50--120 & 0.995 & 4 & 4--6 \\
3D cavity transmon & 100--300 & 0.999 & 0.02 & 6--8 \\
Fluxonium & 200--500 & 0.998 & 0.02 & 0.1--1 \\
\bottomrule
\end{tabular}
}
\label{tab:hardware_compare}
\end{table}

\begin{table}[h]
\centering
\caption{Representative quantum algorithm and hybrid model metrics~\cite{Schuld2022,SchuldPetruccione2021,Nation2024}.}
\resizebox{\textwidth}{!}{
\begin{tabular}{lcccc}
\toprule
Algorithm & Qubits & Circuit Depth & Accuracy (rel.) & Typical Hardware \\
\midrule
QFT / Phase Estimation & $n+1$ & $O(n^2)$ & High (noisy) & Simulators, NISQ \\
Grover Search & $O(\sqrt{N})$ & $O(\sqrt{N})$ & Probabilistic & NISQ / Fault-tolerant \\
VQE & 10--100 & 50--300 & $<1\%$ energy error & Superconducting / Trapped-ion \\
QAOA & 10--500 & $p\times$2 layers & Problem-dependent & Superconducting / Neutral-atom \\
QCNN (this work) & 4--16 & 10--40 & 95--98\% classification & IBM / IQM / simulators \\
QGAN  & 4--12 & 20--60 & 90--96\% fidelity & Hybrid classical/quantum \\
\bottomrule
\end{tabular}
}
\label{tab:algorithm_compare}
\end{table}

\begin{table}[h]
\centering
\caption{Comparison among quantum sensing platforms~\cite{Degen2017,Ezratty2024}.}
\resizebox{\textwidth}{!}{
\begin{tabular}{lccc}
\toprule
Platform & Sensitivity & Bandwidth & Operating Temp. \\
\midrule
SQUID & $<1$ fT$/\sqrt{\mathrm{Hz}}$ & kHz--MHz & 4 K \\
NV centre (diamond) & $10$ nT$/\sqrt{\mathrm{Hz}}$ & MHz & 300 K \\
%Rydberg atom (sensor) & $<1~\mu$V/cm$/\sqrt{\mathrm{Hz}}$ & Hz--THz & 300 K \\
Rydberg atom (photon detection) & 300 K & $\sim 10^{-22}$ W & Hz-THz \\
Rydberg atom (field sensing) & 300~K & $<1~\mu$V/cm/$\sqrt{Hz}$ & Hz--THz \\
Trapped ions & nT$/\sqrt{\mathrm{Hz}}$ & Hz--kHz & 4 K \\
Optomechanical cavity & fN$/\sqrt{\mathrm{Hz}}$ & kHz--MHz & 300 K \\
\bottomrule
\end{tabular}
}
\label{tab:sensing_compare}
\end{table}





%\section{Scalability and Technology Readiness}
Quantum technologies can be evaluated through Technology Readiness Levels (TRLs). Following Ezratty~\cite{Ezratty2024}, Table~\ref{tab:trl_compare} summarizes approximate TRL estimates for major quantum domains.

\begin{table}[h!]
\centering
\caption{Estimated TRL and scalability outlook of major quantum technologies.}
\resizebox{\textwidth}{!}{
\begin{tabular}{lccc}
\toprule
Domain & TRL (2025) & Scalability & Industrial Adoption \\
\midrule
Quantum sensing (NV, Rydberg) & 6--8 & Medium--High & Emerging (metrology, RF sensing) \\
Superconducting hardware & 5--7 & High (chip-based) & Established (IBM, Google, IQM) \\
Trapped-ion hardware & 5--6 & Medium (limited parallelism) & Startups (IonQ, Quantinuum) \\
Photonic processors & 4--6 & Medium & Developing (PsiQuantum, Xanadu) \\
Quantum algorithms (VQE/QAOA) & 4--6 & Medium & Proof-of-concept (Academia/industry collaborations) \\
Quantum ML (QCNN/QGAN) & 3--5 & Low--Medium & Early-stage (research labs) \\
\bottomrule
\end{tabular}
}
\label{tab:trl_compare}
\end{table}

%\section{Energy, Cost, and Practical Considerations}
Operating temperature and cryogenic infrastructure represent dominant cost factors for superconducting systems. Moving from 20 mK to 4 K operation (as targeted by this work) drastically reduces cryogenic power overhead, simplifies integration with classical electronics, and enables compact cryocooler-based systems. Quantum sensors operating at room temperature (e.g., Rydberg and NV) further minimize operational complexity. Hybrid architectures can exploit this diversity: a 4 K transmon processor coupled to a 300 K Rydberg sensor via optical or RF links could enable room-temperature quantum interfaces.

\section{Hybrid Integration and Outlook}
The convergence of sensing, computation, and communication technologies suggests a future of integrated quantum systems. Rydberg sensors may act as quantum transducers or readout interfaces for superconducting qubits, providing scalable RF-to-optical conversion at minimal cryogenic load. Variational algorithms and QML pipelines will serve as digital control and data-fusion layers for these hybrid architectures.

\begin{comment}
\section{Summary}
This comparative analysis integrated metrics from quantum sensing, superconducting hardware, and quantum algorithms. Quantitative tables demonstrated complementary trade-offs: high sensitivity and bandwidth for sensing, strong controllability and coherence for hardware, and algorithmic versatility for hybrid quantum–classical computing. These perspectives collectively delineate a roadmap toward scalable and thermally relaxed quantum systems where sensing and computation coexist in a unified hardware platform.
\end{comment}
